% ------------------------------------------------------------------------------
% Appendix B — drafted Phase C from stage2_gate_report.md (markup numerics + trim)
% ------------------------------------------------------------------------------
\section{Markup Cross-Check}\label{app:markup}

The literature's standard conduct object is the markup of price over marginal
cost \citep{Wolfram1999_duopoly, BorensteinBushnellWolak2002_california,
HortacsuPuller2008_ercot}. This paper's conduct object is different by
design: the strategic margin here is how much capacity is made available at
all, not how offered capacity is priced against energy-market competition.
This appendix reports the markup lens on the same data, and why it is a
cross-check rather than a co-primary outcome.

An implied markup constructed from the residual-demand slope
($-RD/(P \cdot \partial RD/\partial P)$, the inverse-semi-elasticity form)
is numerically ill-behaved exactly where this market lives: the slope has a
mass point at zero (rivals saturated) and a long near-zero shoulder, where the
implied markup explodes --- 37.5\% of finite values exceed 100 in absolute
value, and untrimmed correlations are degenerate. Trimmed to
$|\text{markup}| \leq 10$ --- a generous band that still excludes 63--67\% of
intervals, reported rather than hidden --- the markup's correlation with
essentiality is $+0.076$ for Torrens and $+0.041$ pooled: essential intervals
carry, if anything, slightly \emph{higher} implied competitive pressure, the
same wrong-sign-for-market-power pattern as the slope itself
(Section~\ref{sec:results}).

The construction's fragility is the substantive point. A markup lens presumes
the unit is participating in the energy market and asks how hard it presses
its advantage; a unit that declares zero availability in half of all
intervals is invisible to that lens. For conduct that operates on the
availability margin, the markup benchmark is not wrong but beside the point
--- which is why the cheap-capacity share, not the markup, carries this
paper's analysis.
